% 建模活动单 · 利润最大化
% 两个活动：二次函数最值（售价—销量线性关系）与分段函数（门票两段定价）。
\section{活动十四：售价与销量的线性关系——利润最大定价}

\subsection{情境}
某店销售一种商品：每件售价 60 元时，每天可卖出 300 件。
销售记录表明，售价每上涨 1 元，日销量减少 10 件；每降价 1 元，日销量增加 10 件。
每件成本 40 元，固定成本不计。

\subsection{提出问题}
售价定为多少元时，每天的利润最大？最大利润是多少？
“单价越高越好”这个直觉对吗？

\subsection{选模型}
设售价为 $p$ 元/件，日销量为 $q$ 件。
由题意，$q$ 是 $p$ 的一次函数：
\[
  q=q_0-k(p-p_0)=300-10(p-60)=900-10p.
\]
利润 = 单件利润 $\times$ 销量：
\[
  L(p)=(p-40)q=(p-40)(900-10p)=-10p^2+1300p-36000.
\]
$L$ 是开口向下的二次函数，故有最大值。

\subsection{求解：两种读法}
\textbf{读法一（顶点公式）}：
\[
  p^*=-\frac{1300}{2\times(-10)}=65.
\]
代回检验：$q(65)=900-10\times65=250$（件），单件利润 $65-40=25$（元），
\[
  L_{\max}=25\times250=6250\text{（元）}.
\]

\textbf{读法二（导数）}：
\[
  L'(p)=-20p+1300,\qquad L'(p)=0\implies p=65,
\]
且 $L''(p)=-20<0$，故 $p=65$ 是极大值点，也是最大值点。
导数符号还给出趋势：$p<65$ 时 $L'>0$（提价能增收）；
$p>65$ 时 $L'<0$（提价反而减收）。

若已学过均值不等式，还有第三种读法：
\[
  L(p)=10(p-40)(90-p),\qquad (p-40)+(90-p)=50\text{ 为定值},
\]
故当 $p-40=90-p$，即 $p=65$ 时乘积最大，
$L\le10\times25^2=6250$（元）。三种读法指向同一个答案。

\subsection{检验边界}
销量非负要求 $q\ge0$，即 $p\le90$；不亏本要求 $p\ge40$。
所以实际定义域是 $[40,90]$，而顶点 $p=65\in[40,90]$，故确为最大值。
端点处 $L(40)=L(90)=0$。列表检验：
\begin{center}
\footnotesize
\setlength{\tabcolsep}{8pt}
\renewcommand{\arraystretch}{1.3}
\begin{tabular}{cccc}
\toprule
售价 $p$（元） & 日销量 $q$（件） & 单件利润（元） & 日利润 $L$（元） \\
\midrule
40 & 500 & 0 & 0 \\
50 & 400 & 10 & 4000 \\
60 & 300 & 20 & 6000 \\
65 & 250 & 25 & 6250 \\
70 & 200 & 30 & 6000 \\
80 & 100 & 40 & 4000 \\
90 & 0 & 50 & 0 \\
\bottomrule
\end{tabular}
\end{center}

\subsection{检验与反思}
从 60 元提到 65 元，销量从 300 件降到 250 件，利润却从 6000 元升到 6250 元；
再提到 70 元，利润又回落到 6000 元。“多卖”与“贵卖”各有胜负：
$60\to65$ 是提高单价获胜，$65$ 以上是销量下滑过快拖垮利润。
\key{利润最大的价格不在端点，也不在“卖得最多”或“单价最高”处，而在边际利润为零处。}

\section{活动十五：门票的两段定价——分段函数}

\subsection{情境}
某博物馆散客票 60 元/人；同一旅行社组织团队，人数满 60 人可购团体票，
每人 48 元（八折）。历年数据显示：票价每降低 1 元，日均参观人数增加 2 人；
票价为 60 元时，日均人数为 180 人。

\subsection{提出问题}
票价定为多少元时，门票收入最高？
“团体票更便宜，收入就一定变少”对吗？

\subsection{选模型}
设票价为 $p$ 元/人，日均人数为 $n$，则
\[
  n=180+2(60-p)=300-2p.
\]
能否购买团体票取决于人数是否达到 60：$n\ge60\iff p\le120$。
于是收入是一个分段函数
\[
  R(p)=\begin{cases}
    0.8\,p\,(300-2p), & p\le120,\\[2mm]
    p\,(300-2p), & p>120,
  \end{cases}
\]
其中 $p>120$ 时人数不足 60，只能按散客全价售票。

\subsection{求解}
第一段：$R=0.8(-2p^2+300p)$，与 $-2p^2+300p$ 同对称轴
\[
  p=\frac{300}{4}=75,
\]
且 $75\le120$，故该段最大值在 $p=75$ 处取得：
\[
  R(75)=0.8\times75\times150=9000\text{（元）},\quad n(75)=150\text{（人）}.
\]
第二段：$R=p(300-2p)$，对称轴仍是 $p=75$，但定义域为 $p>120$，
在其上单调递减，无极大值。
比较两段，得最优定价 $p^*=75$ 元，最大收入 $9000$ 元。

\subsection{衔接点处的连续性讨论}
衔接点对应 $n=60$，即 $p=120$。分别取左右极限：
\[
  \lim_{p\to120^-}R(p)=0.8\times120\times60=5760,\qquad
  \lim_{p\to120^+}R(p)=120\times60=7200.
\]
左右极限不相等，$R$ 在 $p=120$ 处\key{跳跃间断}，跳跃量 $1440$ 元。
按规则 $n\ge60$ 可买团体票，故 $R(120)=5760$；
但若这 60 人按散客全价购票，收入是 $7200$ 元。
定价时必须先把规则写清楚：同一人数、同一票价，两种售票方式会得到两个不同收入。

\subsection{陷阱：涨价反而收入下降的区间}
在第一段内，$R$ 在 $p=75$ 处取最大值，故 $75<p\le120$ 时 $R$ 单调下降：
\[
  R(90)=0.8\times90\times120=8640,\quad
  R(105)=0.8\times105\times90=7560,\quad
  R(120)=5760.
\]
从 75 元涨到 120 元，收入反而从 9000 元降到 5760 元，减少 $3240$ 元（$36\%$）。
跨过 120 元时虽向上跳跃到 $7200$ 元附近，但 $p>120$ 后仍继续递减，
如 $R(130)=130\times40=5200$（元）。
所以除 $p=120$ 处的跳跃外，$(75,150)$ 内涨价都不能增收。
列表对照：
\begin{center}
\footnotesize
\setlength{\tabcolsep}{6pt}
\renewcommand{\arraystretch}{1.3}
\begin{tabular}{ccccccccc}
\toprule
票价 $p$（元） & 60 & 75 & 90 & 105 & 120 & 121 & 130 & 150 \\
\midrule
人数 $n$ & 180 & 150 & 120 & 90 & 60 & 58 & 40 & 0 \\
收入 $R$（元） & 8640 & 9000 & 8640 & 7560 & 5760 & 7018 & 5200 & 0 \\
\bottomrule
\end{tabular}
\end{center}

\subsection{检验与反思}
两段定价下求最值，必须先确定每一段的定义域与各段的最值，再统一比较；
衔接点处的跳跃会让“凭直觉调价”失灵，甚至出现“涨价反而收入大降”。
模型的局限也要说清楚：$n(p)$ 只是近似的线性关系；
团体票可能带来额外的讲解、预约成本；票价也不能无限压低。
\key{分段函数的最值在各段分别求、最后比较；衔接点的跳跃是不连续处，也是最容易误判的地方。}
